\section{First consequences}
\label{sec:consequences}

Throughout this section $A\in\PB^C_{\gauge_1}$.

\subsection{The centre is classical}

\begin{theorem}[The classical base]\label{thm:base}
The centre $\Zc(A)$ is a closed commutative subalgebra on which \eqref{ax:sq}
and \eqref{ax:sr} hold exactly.  Consequently
\[
  \Zc(A)\;\cong\;C(X_A,\RR)
\]
isometrically and unitally, for a compact Hausdorff space
$X_A=\Spec\Zc(A)$, uniquely determined up to homeomorphism.
\end{theorem}

\begin{proof}
The centre of a Banach algebra is a closed subalgebra containing $1$, and it is
commutative.  For $a\in\Zc(A)$ we have $\gauge_1(a)=0$ by
Lemma~\ref{lem:gauge-basic}(e), so \eqref{ax:pb1} reads
$\norm{a}^2\le\norm{a^2}$, which together with submultiplicativity gives
\eqref{ax:sq}; and \eqref{ax:pb2} reads exactly \eqref{ax:sr}.  Both are
statements about norms of elements of $A$, hence hold in $\Zc(A)$ with the
induced norm.  Now apply Theorem~\ref{thm:gelfand}.
\end{proof}

\begin{definition}\label{def:base}
$X_A$ is the \emph{classical base} of $A$.  An object with $X_A$ a single
point, equivalently $\Zc(A)=\RR1$, is called \emph{pointless}.
\end{definition}

\begin{remark}[Extracted, not imposed]\label{rem:extracted}
The base is read off the algebra; one is not required to sit over an ambient
space chosen in advance.  For $A=M_2(\RR)$ or $A=\HH$ the base is a single
point, so genuinely pointless noncommutative objects are admitted --- the
theory does not secretly reduce to bundles over a nontrivial space.  For
$A=C(X,\RR)$ the base is $X$ and the algebra is its own centre.  The content of
\S\ref{sec:base} is that this assignment is functorial, and \S\ref{sec:bundle}
that it exhibits every object as a bundle of pointless ones.
\end{remark}

\subsection{Approximate spectral reality}

\begin{theorem}[Characters are approximately real]\label{thm:approx-real}
Let $\chi$ be a character of $A$.  Then for every $a\in A$,
\[
  \bigl(\Ima\chi(a)\bigr)^2\;\le\;C\,\gauge_1(a).
\]
In particular every character of $A$ is real-valued on $\Zc(A)$, and if $A$ is
commutative then every character is real-valued.
\end{theorem}

\begin{proof}
Write $\chi(a)=\alpha+i\beta$ and set $b=a-\alpha1$, so $\chi(b^2)=-\beta^2$ and
$\gauge_1(b)=\gauge_1(a)$ by Lemma~\ref{lem:gauge-basic}(b).  As in
Lemma~\ref{lem:sr-characters},
$\bigl\lVert\norm{b^2}1-b^2\bigr\rVert\ge\norm{b^2}+\beta^2$, and
\eqref{ax:pb2} applied to $b$ bounds the left side by
$\norm{b^2}+C\gauge_1(b)$.
\end{proof}

\begin{remark}
This is the quantitative form of ``the theory has real spectra''.  Off the
centre the bound degrades exactly in proportion to the budget, which is what
lets $\HH$ and the rotation matrix $J=\begin{psmallmatrix}0&-1\\1&0\end{psmallmatrix}$
into the theory: both have $J^2=-1$, an unambiguously non-real spectral
direction, and both are funded by $\gauge_1(J)=4$.  The commutative algebra
$\CC$, which has the same spectral defect and \emph{no} budget, is excluded.
The slogan is that complex directions are allowed only when they are
noncommutative ones.
\end{remark}

\subsection{The phase transition at $C=\tfrac14$}

\begin{theorem}[Quantitative normaloid bound]\label{thm:normaloid}
Let $A\in\PB^C_{\gauge_1}$ with $C<\tfrac14$.  Then for every $a\in A$,
\[
  r(a)\;\ge\;(1-4C)\,\norm{a}.
\]
\end{theorem}

\begin{proof}
By \eqref{ax:pb1} and the universal bound $\gauge_1(x)\le4\norm{x}^2$,
\[
  \norm{x^2}\;\ge\;\norm{x}^2-C\gauge_1(x)\;\ge\;(1-4C)\norm{x}^2
  \qquad\text{for every } x\in A .
\]
Apply this with $x=a^{2^n}$ and put $u_n=\norm{a^{2^n}}^{1/2^n}$, so that
$u_{n+1}^{2^{n+1}}\ge(1-4C)\,u_n^{2^{n+1}}$, i.e.\
$u_{n+1}\ge(1-4C)^{1/2^{n+1}}u_n$.  Iterating from $u_0=\norm{a}$,
\[
  u_n\;\ge\;\norm{a}\prod_{k=1}^{n}(1-4C)^{2^{-k}}
  \;=\;\norm{a}\,(1-4C)^{\,1-2^{-n}}\;\ge\;\norm{a}\,(1-4C),
\]
and $r(a)=\lim_n u_n$ by Beurling's formula.
\end{proof}

\begin{corollary}\label{cor:subcritical-rigid}
If $C<\tfrac14$ then every $A\in\PB^C_{\gauge_1}$ is semisimple and contains no
nonzero nilpotent and no nonzero quasinilpotent element.
\end{corollary}

\begin{proof}
A quasinilpotent has $r(a)=0$, hence $\norm{a}=0$ by
Theorem~\ref{thm:normaloid}; nilpotents are quasinilpotent, and $\rad(A)$
consists of quasinilpotents.
\end{proof}

\begin{theorem}[Phase transition]\label{thm:phase-transition}
For the gauge $\gauge_1$:
\begin{enumerate}
\item[(a)] For every $C<\tfrac14$, the class $\PB^C_{\gauge_1}$ consists of
  semisimple algebras with $r\ge(1-4C)\norm{\cdot}$, and contains no algebra
  with a nonzero nilpotent --- in particular neither $M_2(\RR)$ nor $T_2(\RR)$.
\item[(b)] At $C=\tfrac14$ the pointless objects $M_2(\RR)$ and $\HH$ enter,
  with $\kap_{\gauge_1}(M_2(\RR))=\lam_{\gauge_1}(M_2(\RR))=\tfrac14$ attained
  at $E_{12}$ and at $J$ respectively.
\end{enumerate}
Part (a) is \statusproved; the equalities in (b) are \statusnum, the
inequalities $\ge\tfrac14$ being Examples~\ref{ex:H} and \ref{ex:M2}.
\end{theorem}

\begin{proof}[Proof of (a)]
The first clause is Theorem~\ref{thm:normaloid} and
Corollary~\ref{cor:subcritical-rigid}.  Both $M_2(\RR)$ and $T_2(\RR)$ contain
the nonzero nilpotent $E_{12}$, so neither is subcritical.
\end{proof}

\begin{theorem}[A constant \emph{does} separate $T_2$ from $M_2$]
\label{thm:constant-separates}
$\kap_{\gauge_1}(T_2(\RR))\ge\kap^*=0.3098421747\ldots>\tfrac14$, where
$\kap^*$ is the algebraic number of Proposition~\ref{prop:T2-constant}.  Hence
$T_2(\RR)\notin\PB^{1/4}_{\gauge_1}$, while $M_2(\RR)\in\PB^{1/4}_{\gauge_1}$,
and every constant $C\in[\tfrac14,\kap^*)$ admits the one and excludes the
other.
\end{theorem}

\begin{proof}
Proposition~\ref{prop:T2-constant} and Example~\ref{ex:M2}.
\end{proof}

\begin{remark}[Correction to the earlier notes]\label{rem:separation-correction}
Theorem~\ref{thm:constant-separates} contradicts the assertion of the earlier
notes that no $\gauge_1$-constant separates $T_2(\RR)$ from $M_2(\RR)$, an
assertion that rested on locating the extremiser of $T_2(\RR)$ at the shared
element $E_{12}$.  It does not sit there; see
Remark~\ref{rem:T2-correction}.  The consequences are structural, not
bookkeeping.  The earlier notes present the choice between $\gauge_1$ and
$\gauge_\infty$ as a forced trade --- $\gauge_1$ is closure-friendly and makes
the base functorial but cannot exclude non-central radical, $\gauge_\infty$
excludes it but vanishes on whole directions and so breaks closure.  At
$C=\tfrac14$ the trade is not forced in the one case we can compute: the norm
gauge excludes $T_2(\RR)$ \emph{and} keeps the functorial base of
\S\ref{sec:base}.  How general that is, we do not know, and it is the most
interesting question the correction opens up.
\end{remark}

\begin{conjecture}[Semisimplicity at the critical constant]
\label{conj:semisimple}
Every $A\in\PB^{1/4}_{\gauge_1}$ is semisimple, i.e.\ $\rad(A)=0$.
\end{conjecture}

\begin{remark}[Evidence]\label{rem:semisimple-evidence}
Note first that Conjecture~\ref{conj:semisimple} does \emph{not} follow from
Theorem~\ref{thm:normaloid}, which is vacuous at $C=\tfrac14$, and that it is
not a statement about nilpotent \emph{elements}: $M_2(\RR)\in\PB^{1/4}$ contains
the nilpotent $E_{12}$ and is nonetheless simple.

The evidence is that every elementary radical-carrying algebra we can compute is
excluded at $\tfrac14$.  For $\RR[\varepsilon]/(\varepsilon^2)$ the defect is
unfunded outright (Example~\ref{ex:excluded}).  For the others it is a
quantitative matter, and to exclude an algebra it suffices to exhibit one
element with ratio above $\tfrac14$ --- no supremum is needed.  Such witnesses
are found by \texttt{verify/experiments/radical\_scan.py}:
\begin{center}
\begin{tabular}{llc}
\toprule
Algebra & witness & $\bigl(\norm{a}^2-\norm{a^2}\bigr)/\gauge_1(a)$ \\
\midrule
$T_2(\RR)$   & $x_{0.45}$ & $0.30894$ \\
$T_3(\RR)$   & $x_{0.45}$ in the $\{1,3\}$ corner & $0.30894$ \\
$P_{1,2}(\RR)$ & $x_{0.45}$ in the $\{1,2\}$ corner & $0.30190$ \\
$P_{2,1}(\RR)$ & $x_{0.45}$ in the $\{1,3\}$ corner & $0.30047$ \\
\bottomrule
\end{tabular}
\end{center}
Here $x_c=\begin{psmallmatrix}c&1-c^2\\0&-c\end{psmallmatrix}$ is the traceless
family of Proposition~\ref{prop:T2-constant} and $P_{1,2}(\RR)$, $P_{2,1}(\RR)$
are the block upper-triangular (parabolic) subalgebras of $M_3(\RR)$.  The
parabolic cases are the informative ones: unlike $T_n(\RR)$ they contain a full
$M_2(\RR)$ block on the diagonal, so they have many more test elements with
which to fund a radical defect --- and they still cannot fund it at $\tfrac14$.
All four entries are \statusnum; each is stable to $10^{-4}$ under a tripling of
the optimisation effort, and the first is corroborated by the exact computation
of Proposition~\ref{prop:T2-constant}.

If the conjecture holds, then at the critical constant the norm gauge achieves
without an involution what the $C^*$-identity achieves with one --- the
exclusion of the radical --- while retaining the functoriality that
$\gauge_\infty$ and every $\gauge_q$, $q\ge2$, destroy
(Warning~\ref{warn:vanishing}), and the ``discrimination versus closure''
trade-off of Remark~\ref{rem:gauge-tradeoff} dissolves at exactly one value of
$C$.  This is Problem~\ref{prob:semisimple}.
\end{remark}

\begin{remark}[What the transition means]
The interval $C\in(0,\tfrac14)$ is not a family of interesting theories: by
Corollary~\ref{cor:subcritical-rigid} every object there is semisimple and
nilpotent-free, and the surviving finite-dimensional objects are products of
copies of $\RR$ and $\HH$ (Proposition~\ref{prop:strict-survivors}).  The
noncommutative content appears at $C=\tfrac14$, where the pointless objects
$M_2(\RR)$ and $\HH$ arrive together.  Above $\kap^*$ the non-central radical
arrives too.  So the constant is not a free parameter to be tuned away: it
indexes a genuine filtration, with at least two thresholds,
$\tfrac14$ and $\kap^*$, and the interval between them is where the theory looks
most like a noncommutative geometry --- noncommutative points, no radical, and a
functorial classical base.
\end{remark}

\begin{proposition}[The survivors at zero budget]\label{prop:strict-survivors}
Let $A$ be finite-dimensional and satisfy \eqref{ax:sq} exactly together with
\eqref{ax:pb2}.  Then $A$ is a product of copies of $\RR$ and $\HH$.
\end{proposition}

\begin{proof}
\eqref{ax:sq} gives $\rad(A)=0$ (Lemma~\ref{lem:sq-radius}), so by
Wedderburn $A\cong\prod_iM_{n_i}(D_i)$ with $D_i\in\{\RR,\CC,\HH\}$.  A block
with $n_i\ge2$ contains a nonzero square-zero element, contradicting
\eqref{ax:sq}.  A block $\CC$ is central, so its budget vanishes and
\eqref{ax:pb2} reduces to \eqref{ax:sr}, which $i\in\CC$ violates as in
Remark~\ref{rem:fd-shadow}.  Blocks $\RR$ and $\HH$ both survive: the norm of
$\HH$ is multiplicative, giving \eqref{ax:sq}, and
$\gauge_1(v)=4\abs{v}^2$ for $v$ purely imaginary
(Proposition~\ref{prop:real-stampfli}) funds \eqref{ax:pb2} at $C=\tfrac14$
with equality.
\end{proof}

\begin{remark}[Dyson's threefold way, with the unitary class excised]
Proposition~\ref{prop:strict-survivors} says the zero-budget theory sees the
orthogonal and symplectic classes and not the unitary one.  The
infinite-dimensional analogue --- that \eqref{ax:sq} together with
\eqref{ax:pb2} forces an isometric embedding into an algebra of sections with
$\RR$- and $\HH$-fibres --- would be a quaternionic Stone--Weierstrass theorem.
It is \statusconj; see Problem~\ref{prob:quaternionic-sw}.
\end{remark}
